\chapter{线性数据结构：数组、字符串、链表、栈和队列}

数据结构不是算法之前的一张 API 表，而是算法能够成立的物理基础。刷题时遇到一个问题，第一步不是急着套模板，而是判断答案依赖什么信息：需要按位置访问、按最近关系回退、按先进先出扩展，还是只在两端维护候选。数组、字符串、链表、栈、队列和双端队列构成了这条主线。它们的抽象接口很简单，但背后的存储方式、缓存局部性、迭代器失效规则和边界条件，决定了暴力方法能不能被优化成面试可接受的复杂度。

线性结构的共同特点是元素被组织成一条序列。区别在于这条序列是否连续存储，是否支持下标随机访问，插入删除发生在什么位置，算法状态如何随着左右边界移动而更新。真正掌握这一章，意味着看到题目时可以回答四个问题：序列是否固定，窗口是否连续，历史状态能否增量维护，删除或回退是否只发生在某一端。

从能力边界看，\code{vector} 负责连续存储、随机访问和顺序扫描，是 DP 表、排序数组、邻接表和双指针题的默认选择；\code{string} 负责字符序列和子串窗口，但要警惕 \code{substr} 的复制成本；链表负责节点重连，适合反转、合并、环检测和快慢指针；\code{stack} 负责最近未解决状态，常用于括号、表达式、路径和单调栈；\code{queue} 负责按层扩展，是 BFS 的自然结构；\code{deque} 负责两端维护候选，尤其适合滑动窗口极值。把这些结构先放在同一张能力地图里，后面每个算法就不会变成孤立模板。

数组和 \code{std::vector} 是刷题默认容器。连续存储带来两个关键优势：下标访问是 \complexity{O(1)}，相邻元素通常落在相近内存位置，CPU 缓存命中率高。很多初学者只记住第一个优势，忽略第二个优势。实际程序中，同样是线性扫描，\code{vector} 往往比链表快得多，并不是因为复杂度写法不同，而是因为连续内存更容易被硬件预取。代价也来自连续存储：中间插入和删除必须移动后续元素；容量不足时扩容会重新申请一段更大的空间，把旧元素移动或复制过去，原来的指针、引用和迭代器都可能失效。

\begin{lstlisting}[language=C++,caption={vector 的容量、大小和失效规则}]
std::vector<int> nums;
nums.reserve(n);           // 只预留容量，不构造元素，size 仍然是 0。
nums.push_back(10);
nums.emplace_back(20);

std::vector<int> dp(n + 1, 0);  // 直接创建 n + 1 个元素。

const int* first = nums.data();
nums.push_back(30);        // 如果触发扩容，first 立刻失效。
\end{lstlisting}

\code{reserve} 和 \code{resize} 必须分清。\code{reserve} 只改变容量，不构造元素，适合已经知道将要追加多少元素的场景；\code{resize} 改变大小，会构造或销毁元素，适合直接创建 DP 数组、计数数组或矩阵行。写 \code{vector<int> dp(n + 1)} 后，\code{dp[i]} 是合法的；只写 \code{reserve(n + 1)} 后访问 \code{dp[i]} 是越界。面试里这类错误很隐蔽，因为代码可能在小样例上运行，到了评测环境才崩溃。

数组题的暴力方法通常是枚举一个区间、一个配对或一个分割点。优化的关键，是让“下一步答案”复用“上一步状态”。例如区间和暴力做法是每次重新累加 \code{nums[left..right]}，复杂度 \complexity{O(n^3)} 或 \complexity{O(n^2)}；前缀和把任意区间和转成两个端点差值，使查询变成 \complexity{O(1)}。滑动窗口则在区间右端加入一个元素、左端删除一个元素，用增量更新维护窗口状态。它们都依赖数组按下标稳定访问。

\begin{lstlisting}[language=C++,caption={从重新求和到前缀和查询}]
std::vector<long long> build_prefix_sum(const std::vector<int>& nums)
{
    std::vector<long long> prefix(nums.size() + 1, 0);
    for (std::size_t i = 0; i < nums.size(); ++i) {
        prefix[i + 1] = prefix[i] + nums[i];
    }
    return prefix;
}

long long range_sum(const std::vector<long long>& prefix,
                    std::size_t left,
                    std::size_t right)
{
    return prefix[right + 1] - prefix[left];
}
\end{lstlisting}

这段代码里 \code{prefix[i]} 表示前 \code{i} 个元素之和，因此区间 \code{[left, right]} 的和是 \code{prefix[right + 1] - prefix[left]}。多出来的第 0 位看似绕，但它消除了 \code{left == 0} 的特殊判断。很多高质量题解都会主动设计这种哨兵位，因为它能让边界条件服从同一条公式。

字符串 \code{std::string} 可以看成字符数组，但不能只按字符数组理解。第一，\code{substr} 会创建新字符串，频繁调用会引入大量复制；第二，\code{size()} 返回字节数，不等于所有自然语言里的“字符数”；第三，题目如果限定小写英文，计数结构可以用 \code{std::array<int, 26>}，如果字符集不固定，就应使用哈希表。刷题常见输入大多是 ASCII 或小写英文，面试时仍然要主动说明自己的假设。

\begin{lstlisting}[language=C++,caption={异位词窗口：避免反复 substr}]
std::vector<int> find_anagrams(const std::string& s, const std::string& p)
{
    constexpr int ALPHABET_SIZE = 26;
    std::vector<int> answer;
    if (s.size() < p.size()) {
        return answer;
    }

    std::array<int, ALPHABET_SIZE> need{};
    std::array<int, ALPHABET_SIZE> window{};
    for (const char ch : p) {
        ++need[ch - 'a'];
    }

    for (std::size_t right = 0; right < s.size(); ++right) {
        ++window[s[right] - 'a'];
        if (right >= p.size()) {
            --window[s[right - p.size()] - 'a'];
        }
        if (window == need) {
            answer.push_back(static_cast<int>(right + 1 - p.size()));
        }
    }
    return answer;
}
\end{lstlisting}

暴力版本会枚举每个长度为 \code{p.size()} 的子串，复制出来再排序或统计。复制子串本身就是 \complexity{O(k)}，排序还要 \complexity{O(k\log k)}。窗口版本只维护 26 个计数，右端进入一个字符，左端离开一个字符，每步常数更新。优化成立的原因不是“用了滑动窗口”这个名词，而是题目只关心固定长度窗口内的字符频次，窗口相邻状态只差两个字符。

链表是线性结构里最容易被误解的一个。教材常说链表插入删除是 \complexity{O(1)}，但这句话有前提：你已经拿到了要操作位置的前驱节点或当前节点。如果还需要从头查找位置，查找本身就是 \complexity{O(n)}。链表的价值不在于比数组更快，而在于节点位置稳定、局部重连便宜、结构可以自然表达“下一个”。面试链表题考的不是容器 API，而是指针不丢、边界不漏、连接顺序清楚。

\begin{lstlisting}[language=C++,caption={迭代反转链表：先保存后继，再改变指向}]
struct ListNode {
    int value = 0;
    ListNode* next = nullptr;
};

ListNode* reverse_list(ListNode* head)
{
    ListNode* previous = nullptr;
    ListNode* current = head;

    while (current != nullptr) {
        ListNode* next = current->next;
        current->next = previous;
        previous = current;
        current = next;
    }

    return previous;
}
\end{lstlisting}

这段代码只有四行核心操作，顺序却不能乱。先保存 \code{next}，因为一旦执行 \code{current->next = previous}，原来的后继就找不到了；再把当前节点接到已反转部分前面；最后整体向后推进。写链表题时，建议在纸上画三个指针：已经处理好的部分、当前节点、尚未处理的部分。只要这三段关系清楚，反转、两两交换、K 个一组反转都可以从同一套指针动作推出来。

虚拟头节点是链表题的另一个重要技巧。删除倒数第 N 个节点、合并两个链表、分隔链表时，真实头节点可能被删除或替换。如果每次都单独判断头节点，代码会被边界分支切碎。虚拟头节点把“修改头节点”和“修改普通节点”统一成“修改某个前驱节点的 \code{next}”。

\begin{lstlisting}[language=C++,caption={虚拟头节点统一删除逻辑}]
ListNode* remove_nth_from_end(ListNode* head, int n)
{
    ListNode dummy;
    dummy.next = head;

    ListNode* fast = &dummy;
    ListNode* slow = &dummy;

    for (int step = 0; step < n; ++step) {
        fast = fast->next;
    }
    while (fast->next != nullptr) {
        fast = fast->next;
        slow = slow->next;
    }

    ListNode* removed = slow->next;
    slow->next = removed->next;
    return dummy.next;
}
\end{lstlisting}

快慢指针的本质是让两个指针之间保持固定距离或固定速度差。删除倒数第 N 个节点时，快指针先走 N 步，之后快慢一起走；当快指针到尾部，慢指针正好停在待删节点前驱。环检测时，快指针每次走两步，慢指针每次走一步；如果存在环，速度差会让它们在环内相遇。不要把快慢指针背成模板，应该说清楚“两个指针之间的不变量是什么”。

栈是后进先出结构，适合处理“最近出现、最先解决”的关系。括号匹配中，最近的左括号必须先被右括号匹配；路径简化中，最近进入的目录会被 \code{..} 回退；表达式求值中，最近的运算符和操作数构成局部归约。栈的强大之处在于把嵌套结构压平成一条处理过程。

\begin{lstlisting}[language=C++,caption={有效括号：最近未匹配结构}]
bool is_valid_parentheses(const std::string& s)
{
    std::vector<char> stack;
    stack.reserve(s.size());

    for (const char ch : s) {
        if (ch == '(' || ch == '[' || ch == '{') {
            stack.push_back(ch);
            continue;
        }
        if (stack.empty()) {
            return false;
        }

        const char left = stack.back();
        stack.pop_back();
        const bool matched = (left == '(' && ch == ')') ||
                             (left == '[' && ch == ']') ||
                             (left == '{' && ch == '}');
        if (!matched) {
            return false;
        }
    }

    return stack.empty();
}
\end{lstlisting}

这里直接使用 \code{std::vector<char>} 作为栈，比 \code{std::stack<char>} 更方便查看和预留容量。\code{std::stack} 是容器适配器，只暴露 \code{push}、\code{pop}、\code{top}、\code{empty} 等接口；\code{pop()} 不返回元素，因此必须先 \code{top()} 再 \code{pop()}。面试写代码时两者都可以，但要明白适配器牺牲了一部分访问能力，换来更明确的抽象边界。

单调栈是栈在刷题中的高频变体。它维护的不是时间顺序本身，而是一组“尚未找到答案的候选”。每日温度中，栈保存还没有找到更高温度的下标，并且这些下标对应的温度从栈底到栈顶单调不增。新温度到来时，只要它高于栈顶下标的温度，就可以确定栈顶答案；因为这是栈顶下标右侧第一个更高温度。

\begin{lstlisting}[language=C++,caption={每日温度：单调栈保存下标}]
std::vector<int> daily_temperatures(const std::vector<int>& temperatures)
{
    std::vector<int> answer(temperatures.size(), 0);
    std::vector<int> stack;
    stack.reserve(temperatures.size());

    for (int day = 0; day < static_cast<int>(temperatures.size()); ++day) {
        while (!stack.empty() && temperatures[day] > temperatures[stack.back()]) {
            const int previous_day = stack.back();
            stack.pop_back();
            answer[previous_day] = day - previous_day;
        }
        stack.push_back(day);
    }

    return answer;
}
\end{lstlisting}

单调栈常保存下标而不是值，因为题目经常需要距离、宽度，或者需要回到原数组访问更多信息。它能把暴力枚举“右侧第一个更大元素”的 \complexity{O(n^2)} 降到 \complexity{O(n)}，原因是每个下标最多入栈一次、出栈一次。这个摊还分析要会说：虽然某一次循环里可能弹出很多元素，但全局弹出次数不超过 n。

队列是先进先出结构，适合按层扩展。BFS 的正确性来自队列顺序：先入队的是距离更短或时间更早的状态，因此第一次到达一个未访问状态时，就已经得到了最短步数。二叉树层序遍历、腐烂橘子、多源最短扩散、无权图最短路，都依赖这个性质。与栈不同，队列不会深挖某一条路径，而是把同一层的状态先处理完。

\begin{lstlisting}[language=C++,caption={层序遍历：队列保存当前边界}]
std::vector<std::vector<int>> level_order(TreeNode* root)
{
    std::vector<std::vector<int>> levels;
    if (root == nullptr) {
        return levels;
    }

    std::queue<TreeNode*> queue;
    queue.push(root);

    while (!queue.empty()) {
        const int level_size = static_cast<int>(queue.size());
        std::vector<int> level;
        level.reserve(level_size);

        for (int i = 0; i < level_size; ++i) {
            TreeNode* node = queue.front();
            queue.pop();
            level.push_back(node->value);

            if (node->left != nullptr) {
                queue.push(node->left);
            }
            if (node->right != nullptr) {
                queue.push(node->right);
            }
        }
        levels.push_back(std::move(level));
    }

    return levels;
}
\end{lstlisting}

这段代码里的 \code{level_size} 是关键。没有它，也可以 BFS，但很难知道哪些节点属于同一层。层序遍历、最短时间、多源扩散都需要把“当前层的节点数量”固定下来，然后处理完这一层再增加距离或分钟数。

\code{std::deque} 支持两端 \complexity{O(1)} 插入删除，是 \code{std::queue} 的默认底层容器，也是单调队列的常用实现。滑动窗口最大值的暴力方法是每个窗口扫描一遍最大值，复杂度 \complexity{O(nk)}。如果用堆，需要处理过期下标。单调队列更贴合题意：队头永远是当前窗口最大值下标，队尾负责淘汰不可能再成为最大值的候选。

\begin{lstlisting}[language=C++,caption={滑动窗口最大值：单调队列}]
std::vector<int> max_sliding_window(const std::vector<int>& nums, int k)
{
    std::vector<int> answer;
    if (nums.empty() || k <= 0) {
        return answer;
    }
    answer.reserve(nums.size());

    std::deque<int> deque;
    for (int right = 0; right < static_cast<int>(nums.size()); ++right) {
        while (!deque.empty() && nums[deque.back()] <= nums[right]) {
            deque.pop_back();
        }
        deque.push_back(right);

        const int left = right - k + 1;
        if (deque.front() < left) {
            deque.pop_front();
        }
        if (left >= 0) {
            answer.push_back(nums[deque.front()]);
        }
    }

    return answer;
}
\end{lstlisting}

单调队列的优化理由和单调栈相似。一个新元素如果大于等于队尾元素，那么队尾元素在未来任何包含新元素的窗口里都不可能成为最大值，因为它更小且更早过期。于是可以安全删除。每个下标最多进队一次、出队一次，总复杂度 \complexity{O(n)}。

\begin{principlebox}{线性结构选型}
需要稳定下标和顺序扫描，优先 \code{vector}；需要字符窗口，使用 \code{string} 加计数结构，不要频繁 \code{substr}；需要局部重连，考虑链表并画清楚前驱、当前、后继；需要最近未解决状态，用栈；需要按层或最短步数扩展，用队列；需要窗口两端同时维护候选，用双端队列。选择结构时同时说出它让哪一步查询、删除或更新变快。
\end{principlebox}

本章建议配套练习：数组与前缀和做 303、304、560；字符串窗口做 3、438、567、76；链表做 19、21、24、141、142、206、92、25；栈做 20、71、150、155、739、84；队列和双端队列做 102、199、239、994。复盘时不要只写最终代码，要写三行说明：暴力方法慢在哪里，数据结构保存了什么状态，为什么每个元素只被处理有限次。
